\scanpage{122}
\mspage{37}
% archive sheet number: 121

% L01-02
Ce revêtement est de la forme
\[
U'=U(\infty)\mathbin{\times^G}\Gamma,
\tag*{(}
\]

% L03-05
pour un hom $\unsure{\varphi|\phi}:G\to\Gamma$ dont le calcul est immédiat, et qui \deleted{\unsure{se}} s'explicite, moyennant (4.8.1) et (4.14.5), comme

% L06
\[
\varphi(\lambda_1,\ldots,\lambda_r)=\sum_1^r m_i\lambda_i.
\tag*{(4.14.6)}
\]

% L07-09
On en conclut alors de (4.10.1) et (4.13.1) le calcul de $H^*(U',\Lambda)$. On trouve en particulier, grâce à 4.11 :

% L10-16
\uline{Corollaire} 4.15. \inserted{Dans les notations de 4.14,} Soit $m$ le pgcd des $m_i$, \deleted{Alors} \deleted{\unsure{et}} $m'$ le morceau de $m$ premier à $p$, et $\Gamma'\inserted{=m'\Gamma=m\Gamma}$ le sous-groupe de $\Gamma\simeq\prod_{\ell\ne p}\unsure{\mathbb{Z}|Z}_\ell(1)$ d'indice $m'$. Alors le sous-groupe des éléments de $\Gamma$ qui opèrent trivialement sur $H^*(U',\Lambda)$ (où $\Lambda=\mathbb{Z}/n\mathbb{Z}$, $n>1$, $(n,p)=1$) est égal à $\Gamma'$. \deleted{Pour tout \unsure{ent. premier à}}

% L17
On notera que $\Gamma'$ est \uline{indépendant de} $\Lambda$, \unsure{i.e. de $n$}.

% L18-19
\deleted{4.16. Nous allons mettre \unsure{les résultats sous la forme d'un énoncé} \illegible{1}}

% L20-24
4.16. Nous allons mettre les résultats 4.14, 4.15 sous une forme \deleted{\illegible{1}} plus globale, qui nous sera utile. Pour ceci, supposons donné un schéma \inserted{régulier} $X$, \deleted{\unsure{un}} un diviseur $D$ sur $X$ à croisements normaux, et une section
\[
t\in\Gamma(X,\unsure{\mathcal{O}|\mathscr{O}}_X).
\]

% L25-26
\deleted{On suppose que $\unsure{\pi(t)|V(t)|(t)}$} \deleted{Soit $U=$}
